\section{Introduction}
\label{sec:intro}
% ====================================================================

Legged robots that traverse irregular terrain---stairs, hurdles, gaps---must perceive geometry and maintain tight coordination between vision and control.
Training locomotion policies in simulation with reinforcement learning (RL) and transferring them to real hardware has proven effective for agile skills~\citep{hwangbo2019learning, lee2020learning, rudin2022learning}, and recent work~\citep{zhuang2023robot, cheng2023parkour, hoeller2024anymal, chane2025soloparkour} has extended this to parkour---climbing, jumping, and crawling---via teacher--student distillation.
In these systems, a \emph{teacher} trains with privileged terrain information (e.g.\ ground-truth heightmaps), then a \emph{student} reproduces the teacher's behavior from onboard sensors.

A critical assumption, however, is access to an onboard depth sensor.
Depth cameras (e.g.\ Intel RealSense) transfer well from simulation~\citep{miki2022learning, agarwal2023legged}, but not all platforms include them---the Unitree Go2, for instance, ships with only RGB cameras and a LiDAR sensor, with no onboard depth camera~\citep{unitreego2}.

Deploying RGB-trained policies on real robots requires bridging the \emph{sim-to-real gap}.
Three strategies have emerged: (1)~\textbf{domain randomization}~\citep{tobin2017domain}, which randomizes visual properties; (2)~\textbf{generative augmentation}~\citep{yu2024lucidsim}, which uses diffusion models for photorealistic training images; and (3)~\textbf{canonical representations}, which project both domains into a shared space via a foundation model.

We pursue strategy~(3) using Depth~Anything~V2 (DA2)~\citep{yang2024depth}, a pre-trained monocular depth estimator trained on millions of diverse real images.
We propose that because DA2 maps arbitrary RGB to a depth space that is
largely robust to appearance variation, a student policy operating on
estimated depth may inherit that robustness without photorealistic
rendering or explicit domain randomization.

\paragraph{Thesis statement.}
We hypothesize that pre-trained monocular depth estimation can serve as a
visually robust intermediate representation for vision-based quadruped
parkour, allowing a student trained on estimated depth in simulation to
generalize across visual domains.
We evaluate this hypothesis entirely in simulation; real-world transfer
remains future work.

\paragraph{Contributions.}
\begin{itemize}[nosep,leftmargin=*]
  \item A two-stage teacher--student pipeline from privileged scandot
    heightmaps to monocular depth for quadruped parkour.
  \item A demonstration that monocular depth enables basic parkour (stairs, hurdles, gaps) in simulation.
  \item Identification of depth signal quality---governed by
    model scale and terrain texture---as the key factor limiting performance.
  \item A domain invariance experiment testing DA2-based students under visual perturbations.
\end{itemize}
